\section{Integral homology and geometric Morse counts}
\label{sec:morse}

We first separate two facts used in the construction: integral homology
gives a lower bound for every admissible stabilized Morse function, and
Smale's theorem realizes that bound by an ordinary Morse function on a
simply connected manifold of sufficiently large dimension.  Throughout
this section, put
\[
 r_i(W)=\rank H_i(W;\Z),\qquad
 t_i(W)=d\bigl(\Tor H_i(W;\Z)\bigr).
\]
Thus $\lambda(W)=\sum_i r_i(W)+2\sum_i t_i(W)$, with both sequences
understood to vanish outside the homological range.

\begin{lemma}[Ranks of free integral complexes]\label{morse:rank}
Let $C_*$ be a finite chain complex of finitely generated free abelian
groups.  If $b_i=\rank H_i(C_*)$ and
$t_i=d(\Tor H_i(C_*))$, then
\[
 \rank C_i\geq b_i+t_i+t_{i-1},\qquad
 \sum_i\rank C_i\geq\sum_i b_i+2\sum_i t_i.
\]
The same statements hold for a $\Z/2$-graded complex, with indices
interpreted modulo two.  In either grading, the lower bound on total
rank is attained by a free complex with any prescribed finite
collection of finitely generated homology groups.
\end{lemma}

\begin{proof}
Write $Z_i=\ker\partial_i$ and
$B_i=\operatorname{im}\partial_{i+1}$, which are free abelian.
The exact sequences defining cycles, boundaries, and homology give
\[
 \rank C_i=b_i+\rank B_i+\rank B_{i-1}.
\]
Smith normal form for the inclusion $B_i\subset Z_i$ shows that
$\Tor(Z_i/B_i)$ has at most $\rank B_i$ nontrivial cyclic invariant
factors.  Hence $t_i\leq\rank B_i$, proving both inequalities.
This argument also applies verbatim with cyclic indices.

For attainment, decompose each prescribed homology group into its free
part and nontrivial cyclic invariant factors.  A free generator in
degree $i$ is represented by an isolated copy of $\Z$ in that degree.
A cyclic factor $\Z/n$, $n>1$, in degree $i$ is represented by
\[
 \Z\xrightarrow{\ n\ }\Z
\]
in degrees $i+1$ and $i$, with all other differentials zero.
Take the direct sum of these complexes.  In the cyclic grading the
source and target lie in the two opposite parities, with the reverse
differential zero.  Each cyclic factor therefore costs exactly two
free generators in either grading.
\end{proof}

In particular, $\beta_{\Z}^{\mathrm{cyc}}(W)$ is the minimum total
rank of a free $\Z/2$-graded integral complex whose homology is
$\overline H_{\bar *}(W;\Z)$.  This interpretation concerns algebraic
complexes; it does not assert geometric realization by Morse functions.
The distinction between this minimum and $\lambda(W)$ is that cyclic
grading permits torsion at different integer degrees of the same
parity to share invariant factors.

\begin{proposition}\label{morse:stable}
For every closed smooth manifold $W$, with the exact quadratic
condition at infinity and all stabilizations $k\geq0$ allowed,
\[
 \SM(W)\geq\lambda(W).
\]
\end{proposition}

\begin{proof}
Let $F\colon W\times\R^k\to\R$ be Morse and equal to a fixed
nondegenerate quadratic form $Q$ outside a compact subset of the
product.  We construct a finite free integral complex with one
generator per critical point and homology equal to that of $W$, up
to a uniform grading shift.  For $k=0$ the ordinary Morse complex
suffices.  Suppose $k>0$.

After a linear change of auxiliary coordinates, write
\[
 Q(u,v)=-|u|^2+|v|^2,
 \qquad (u,v)\in\R^{\nu}\times\R^{k-\nu}.
\]
Choose $R_0>0$ such that $F=Q$ for $|(u,v)|\geq R_0$ and all
critical points lie in $W\times B_{R_0}$.  This is possible because
$dQ$ vanishes only at the origin.  Choose
\[
 A>\max_{W\times\overline B_{R_0}}\{|F|,|Q|\},
 \qquad T>R_0,\quad T^2>A,
\]
and set
\[
 K=W\times\overline B_T,\qquad
 K_\pm=K\cap\{F\leq\pm A\}.
\]
The core $W\times\overline B_{R_0}$ is excluded at level $-A$
and included at level $A$, for both $F$ and $Q$.  Outside that core
the two functions agree.  Consequently
\begin{equation}\label{morse:endpoints}
 K_\pm=K\cap\{Q\leq\pm A\}.
\end{equation}

We explain why this compact truncation introduces no extra cells in
the level band $[-A,A]$.  On the side boundary
$W\times\partial B_T$, the restriction of $F$ is $Q$.
Its restricted differential can vanish only when $u=0$ or $v=0$,
so its only possible critical values are $T^2$ and $-T^2$.
Neither is in the band.  Thus a gradient-like vector field for $F$
on the band can be chosen tangent to the side boundary: take a
tangential gradient there, extend across a collar, and combine it
with interior gradient-like fields by a partition of unity.
Away from the critical points all these fields increase $F$ strictly.

The usual Morse attachment argument now applies to this compact
manifold with boundary.  Between critical levels its deformation
flow preserves the side boundary, and each interior critical point
attaches exactly one cell of its Morse index.  This is the standard
attachment theorem \cite[Section~3]{Milnor1963}, with the tangential
field supplying the boundary-preserving deformation.  The endpoint
level hypersurfaces meet the side boundary transversely.  Hence the
endpoint sublevels are compact manifolds with corners, or empty,
and have finite CW type.  Replacing $K_-$ by a finite CW model and
using cellular approximation on the finitely many attaching maps
gives a finite free relative cellular complex $C_*$ satisfying
\begin{equation}\label{morse:relative-complex}
 \sum_i\rank C_i=\#\operatorname{Crit}(F),\qquad
 H_i(C_*)\cong H_i(K_+,K_-;\Z).
\end{equation}
Critical values may first be separated by a small perturbation
supported near the critical points; this leaves the endpoint pair
and all critical-point counts unchanged.

By \eqref{morse:endpoints}, shrinking $v$ to zero preserves both
members of the pair: it decreases $Q$ and the Euclidean norm.
The resulting deformation retraction identifies its relative
homology with that of
\[
 W\times\bigl(D_T^{\nu},
       \{u\in D_T^{\nu}:|u|\geq\sqrt A\}\bigr).
\]
The disk--annulus pair has integral relative homology $\Z$ in
degree $\nu$ and zero otherwise.  When $\nu=0$ it is the pair
$(\{0\},\varnothing)$, with the same assertion in degree zero.
The relative K\"unneth theorem therefore gives
\[
 H_i(C_*)\cong H_{i-\nu}(W;\Z).
\]
This includes positive definite, negative definite, and indefinite
quadratic forms.  Applying Lemma~\ref{morse:rank} to
\eqref{morse:relative-complex} proves
$\#\operatorname{Crit}(F)\geq\lambda(W)$.
Taking the minimum over all admissible functions proves the claim.
\end{proof}

\begin{remark}\label{morse:conventions}
The common convention allowing $\|dF-dQ\|$ to be uniformly bounded
gives the same stable Morse number. Indeed, write $F=Q+h$ and
$\|dh\|\leq K$ in a product metric. Compactness of $W$ gives
$|h(x,z)|\leq K_0+K|z|$, whereas nondegeneracy gives
$\|dQ(z)\|\geq c|z|$ for some $c>0$. If a radial cutoff $\chi_R$
equals one on $B_R$, vanishes outside $B_{2R}$, and has
$\|d\chi_R\|\leq C/R$, then on the intervening annulus
\[
 \|d(\chi_Rh)\|\leq K+C(K_0/R+2K).
\]
For sufficiently large $R$ this is less than $cR$.
Thus $Q+\chi_Rh$ equals $F$ near every critical point, has no new
critical points, and equals $Q$ outside a compact set. The converse
inclusion is immediate. This identifies our convention with the
one used in the stable Morse comparisons cited in the introduction.
\end{remark}

\begin{theorem}[Smale]\label{morse:smale}
Let $W$ be a simply connected closed smooth manifold of real
dimension at least six.  There is an ordinary Morse function
$h\colon W\to\R$ having exactly
\[
 r_i(W)+t_i(W)+t_{i-1}(W)
\]
critical points of index $i$, simultaneously for every $i$.
Consequently
\[
 \SM(W)=\Morse(W)=\lambda(W).
\]
\end{theorem}

\begin{proof}
The existence assertion is the torsion-inclusive form of the
classical geometric Morse theorem of Smale
\cite[Theorem~2.4 and the following paragraph]{Smale1963}; see also
\cite[Theorem~2 and Section~6]{Cerf1962} for its statement and handle
proof.  In these statements the torsion count is the number of
cyclic invariant factors, equivalently the minimum number of
generators of the torsion subgroup.  Thus summing the prescribed
index counts gives $\Morse(W)\leq\lambda(W)$.
An ordinary Morse function is admissible in the definition of
$\SM$ by choosing $k=0$; on the compact space $W$ the exterior
condition is then vacuous.  Proposition~\ref{morse:stable} yields
\[
 \lambda(W)\leq\SM(W)\leq\Morse(W)\leq\lambda(W),
\]
which proves both equalities.
\end{proof}
